\begin{figure}[tbp]
\centering
\begin{tikzpicture}[gclplate,x=1cm,y=1cm]
  \node[gcllabel,draw=GCLBlue!75!black,align=center,text width=2.8cm,minimum height=1.0cm] (a) at (-4.15,0) {peer terminates};
  \node[gcllabel,draw=GCLWarm!75!black,align=center,text width=2.8cm,minimum height=1.0cm] (b) at (0,0) {queue\\$[m]$};
  \node[gcllabel,draw=GCLGreen!75!black,align=center,text width=2.8cm,minimum height=1.0cm] (c) at (4.15,0) {message remains};
  \draw[gclbluearrow] (a.east) -- (b.west);
  \draw[gclwarmarrow] (b.east) -- (c.west);
\end{tikzpicture}
\caption{Plate 17. An orphan message. A queued message can outlive a peer in an ill-disciplined variant. This is a diagnostic example, not a theorem about every asynchronous calculus.}
\label{plate:protocol-17}
\end{figure}
